\documentclass[10pt,a4paper]{article}

% Configurações de idioma e fonte
\usepackage[utf8]{inputenc}
\usepackage[T1]{fontenc}
\usepackage[brazilian]{babel}
\usepackage[margin=0.75in]{geometry}
\usepackage{hyperref}
\usepackage{enumitem}

% Desabilita a numeração de páginas
\pagestyle{empty}

% Formatação personalizada de seções
\usepackage{titlesec}
\titleformat{\section}{\large\bfseries\uppercase}{}{0pt}{}[\titlerule]
\titlespacing{\section}{0pt}{12pt}{8pt}

\begin{document}

% Cabeçalho
\begin{center}
    {\LARGE \textbf{Kevin Liu Rodrigues}}  \\
    \vspace{2pt}
    Brasília, DF, Brasil $|$ +55 19 99732 6791  $|$ \href{mailto:kkevin.liu@gmail.com}{kkevin.liu@gmail.com} $|$ \href{http://hdl.handle.net/1843/39095}{http://hdl.handle.net/1843/39095}
\end{center}

% Resumo
\section{Resumo}
Sou cientista com doutorado e sólida formação quantitativa, sempre buscando
aplicar boas práticas e agregar conhecimento a tudo que faço. Durante o
doutorado, formalizei matematicamente propriedades complexas de sistemas
dinâmicos que até então eram apenas intuitivas para a comunidade internacional,
o que considero um dos resultados mais marcantes da minha trajetória acadêmica.
Desde então, tenho atuado com aprendizado de máquina aplicado e ciência de
dados em contextos práticos variados, consolidando minha experiência como
programador e cientista de dados em projetos de transformação de dados,
modelagem preditiva e classificação.

% Experiência
\section{Experiência}

\textbf{Moray} (2021 -- Presente)  \\
\textit{Cientista de Dados e Engenheiro de Dados} 
\begin{itemize}[noitemsep, topsep=0pt]
    \item Desenvolvimento de modelos preditivos para auxiliar o manejo de safras. 
    \item Atuação com dados espaço-temporais e dados imperfeitos ou ausentes, além de modelos de aprendizado de máquina.
	\item Infra-estrutura de dados e deploy de modelos
\end{itemize}

\vspace{6pt}

\textbf{Rappi Brasil} (2020 -- 2021)  \\
\textit{Cientista de Dados} 
\begin{itemize}[noitemsep, topsep=0pt]
    \item Desenvolvimento de modelos de valor do usuário (LTV) e outras métricas relacionadas ao usuário.
    \item Visualização de grandes volume de dados e dashboards
\end{itemize}

\vspace{6pt}

\textbf{Olivia AI} (2018 -- 2018)  \\
\textit{Cientista de Dados} 
\begin{itemize}[noitemsep, topsep=0pt]
    \item Classificação de transações monetárias para um aplicativo de organização de finanças pessoais.
	\item Desenvolvimento de métricas e processamento de linguagem natural (NLP).
\end{itemize}

\vspace{6pt}

\textbf{Universidade Federal de Minas Gerais} (2015 -- 2016)  \\
\textit{Pesquisador} 
\begin{itemize}[noitemsep, topsep=0pt]
    \item Pesquisa e desenvolvimento sob concessão do professor Gilberto Medeiros Ribeiro para um microscópio de campo próximo de alta potência.
\end{itemize}

\vspace{6pt}

\textbf{TRIU} (2012 -- 2014)  \\
\textit{Professor de Matemática Voluntário} 
\begin{itemize}[noitemsep, topsep=0pt]
    \item Aulas de matemática em cursinho para pessoas de baixa renda que desejam ingressar em universidades competitivas.
\end{itemize}

\vspace{6pt}

\textbf{Universidade de Campinas} (2010 -- 2012)  \\
\textit{Iniciação Científica} 
\begin{itemize}[noitemsep, topsep=0pt]
    \item Desenvolvimento de eletrônica de detecção capacitiva para medição de clusters atômicos em câmara de vácuo.
    \item Desenvolvimento de lentes eletrostáticas para seleção de clusters por massa.
\end{itemize}

% Educação
\section{Educação}

\textbf{Universidade Federal de Minas Gerais / University of Michigan} (2016 -- 2021)  \\
Doutorado (Ph.D.) em Física Estatística e computacional \\
Orientação: Ronald Dickman (UFMG) e Kevin Wood (UMICH)

\vspace{4pt}

\textbf{Universidade de Campinas} (2012 -- 2015)  \\
Mestrado em Física Aplicada 

\vspace{4pt}

\textbf{Universidade de Campinas} (2008 -- 2012)  \\
Bacharelado em Física 

% Competências
\section{Competências}
\begin{itemize}[label={}, leftmargin=0pt, noitemsep]
    \item \textbf{Linguagens de programação:} Python, SQL, C, C++, TCL/TK, JavaScript
    \item \textbf{Ferramentas:} Amazon Web Services (AWS), Softwares CAD, MLflow
    \item \textbf{Idiomas:} Português (Nativo), Inglês (Fluente falado/escrito) 
\end{itemize}

\end{document}
